\section{Pipeline Validation and Structural Corpus Summary}
\label{sec:validation}

To validate the sequential waterfall architecture across the 2021--2025 observation window, we evaluate data retention and fallback resolution across all four tranches prior to downstream topological modeling.

\subsection{Cohort Continuity and Fallback Resolution (Tranches 1--2)}
As summarized in Table~\ref{tab:pipeline_yield}, filtering the 20 quarterly QQQ Form N-PORT filings (\texttt{S000101292}) for domestic U.S. issuers (\texttt{country == 'US'}) and deduplicating by ISIN yields 120 unique CIKs. Our manual longitudinal continuity audit removes 12 issuers that exited or entered the index mid-window due to corporate restructuring or mergers and acquisitions (such as Activision Blizzard following its acquisition by Microsoft Corporation), locking in a balanced panel of $N = 108$ firms ($540$ firm-year Form 10-K filings).

In Tranche~2, programmatic parsing via \texttt{edgartools} successfully isolated Item~1 and Item~1A narratives for 106 of the 108 firms, while custom string-manipulation scripts resolved structural DOM anomalies in the filings of Honeywell International Inc.\ and Intel Corporation to achieve 100\% narrative coverage ($540/540$ filings). For quantitative segment and geographic revenue tables (Items~7 and~8), auditing the initial pool of 2,160 XBRL statement instances revealed significant taxonomy heterogeneity across filers. By cascading unparsed filings through Stage~2 (automated segment/geographic keyword extraction with manual verification) and Stage~3 (native HTML anchor identification via \texttt{str.split}), the pipeline recovers longitudinal segment and geographic revenue structures across the 540 firm-year snapshots.

\begin{table}[htbp]
\centering
\caption{Descriptive Structural Summary of Extracted Micro-Network Entities Across the Static QQQ Cohort ($N=108$ Firms, 2021--2025). \textit{Note: Bracketed placeholders to be populated upon completion of Tranche 3 LLM serialization run.}}
\label{tab:entity_summary}
\small
\begin{tabular}{@{}lllrr@{}}
\toprule
\textbf{Micro-Branch Hub} & \textbf{Node / Attribute Tier} & \textbf{Source Section} & \textbf{Total Extracted} & \textbf{Mean per Firm-Year} \\ 
\midrule
\texttt{Firm Root ($r_i$)} & 10-K Filing Date ($d_{\text{file}}$) \& Market Cap ($M$) & Header / Metadata & 540 & 1.00 \\
\midrule
\texttt{Product/Services} & Operating Segments ($S_i^{(t)}$) & Items 1, 7, 8 & \texttt{[TBD]} & \texttt{[TBD]} \\
 & Segment Revenue Attributes ($w_{\text{rev}}$) & Items 7, 8 & \texttt{[TBD]} & \texttt{[TBD]} \\
 & Granular Product/Service Leaves & Item 1 & \texttt{[TBD]} & \texttt{[TBD]} \\
\midrule
\texttt{Competitor} & Segment-Specific Competitors & Item 1 & \texttt{[TBD]} & \texttt{[TBD]} \\
 & General / Enterprise Competitors (\texttt{Others}) & Item 1 & \texttt{[TBD]} & \texttt{[TBD]} \\
\midrule
\texttt{Risk} & Segment-Specific Operational Risks & Item 1A & \texttt{[TBD]} & \texttt{[TBD]} \\
 & Macroeconomic / External Risks (\texttt{Other}) & Item 1A & \texttt{[TBD]} & \texttt{[TBD]} \\
\midrule
\texttt{Geographic \& Rev.} & Geographic Region Nodes ($\text{Region}_{a\dots n}$) & Items 7, 8 & \texttt{[TBD]} & \texttt{[TBD]} \\
 & Aggregated \texttt{Rest} / \texttt{Other} Revenue Nodes & Items 7, 8 & \texttt{[TBD]} & \texttt{[TBD]} \\
\bottomrule
\end{tabular}
\end{table}

\subsection{Serialized Micro-Network Corpus Statistics (Tranche 3)}
Table~\ref{tab:entity_summary} reports descriptive structural statistics of the nested JSON corpus generated by the Tranche~3 LLM serialization workflow across the 2021--2025 panel. Across the 108 firms, each micro-network $\Gmicro$ decomposes corporate disclosures into operating segments $S_i^{(t)}$ linked to granular product/service leaves, segment-level and general competitors, internal versus external risk factors, and sovereign geographic revenue nodes. This structured serialization preserves intra-firm heterogeneity while providing standardized node and attribute schemas for macro-level projection in Tranche~4.